\section{参数速查与使用限制}
数值对扩展使用 \key{\gradbarsdeclarepairtype} 注册、
\key{\gradpairgraphic} 调用、\key{\gradbarspairvalue} 读取记录；
\key{\gradbarsdeclarepaircolumn} 显式声明原生双值列。
完整签名、缺失规则、能力查询与示例见第~\ref{sec:pair-extensions-cn}~节。

\begingroup\small
\begin{center}
\begin{tabularx}{\linewidth}{@{}>{\centering\arraybackslash}p{31mm}>{\centering\arraybackslash}p{25mm}>{\centering\arraybackslash}X@{}}
\toprule
\multicolumn{1}{c}{\textbf{参数}} & \multicolumn{1}{c}{\textbf{默认值}} & \multicolumn{1}{c}{\textbf{含义与可选值}}\\
\midrule
\key{min} & \key{0} & 范围下限，必须为零或负值。\\
\key{max} & \key{100} & 范围上限，必须为正值。\\
\key{width} & \key{24mm} & 正的轨道宽度，不含外部标签。\\
\key{height} & \key{1.5ex} & 正的轨道高度。\\
\key{theme} & \key{blue} & 11 种条形主题，见主题章节；显式选择会设置负值互补色。\\
\key{label} & \key{outside} & outside、left、inside、none、end、auto。\\
\key{label width} & \key{4.5em} & 外部标签最小槽宽；自动溢出模式下可扩展。\\
\key{label gap} & \key{.6em} & 轨道与外部标签之间的间隔。\\
\key{baseline} & \key{auto} & auto、text、center 或 bottom；控制行内图形基线。\\
\key{baseline shift} & \key{0pt} & 在预设基线上附加垂直偏移。\\
\key{precision} & \key{1} & 0--6 的整数；保留末尾零。\\
\key{value format} & \key{value} & 原始数值 value 或计算占比 percent。\\
\key{unit} & 空 & 原始数值的后缀；占比模式忽略此项。\\
\key{text} & 空 & 覆盖整个标签；空值恢复自动格式。\\
\key{overflow} & \key{clip} & clip 警告并截断，或 error 报错。\\
\key{color} & 未指定 & 指定单色，并切换为纯色填充。\\
\key{fill} & 由主题指定 & gradient 或 solid。\\
\key{left color} & 由主题指定 & 渐变起点，或纯色填充的颜色。\\
\key{right color} & 由主题指定 & 渐变终点。\\
\key{track color} & 由主题指定 & 空轨道颜色。\\
\key{negative fill} & \key{solid} & 未选主题时为纯色；显式主题设置对应模式。\\
\key{negative left/right color} & 空 & 主题可设置负值两端；覆盖任一项会启用渐变。\\
\key{text color} & 由主题指定 & 标签文字颜色；auto 自动选黑白。\\
\key{font} & 沿用周围字体 & 标签字体声明。\\
\key{style} & 未指定 & 应用已定义的命名样式。\\
\key{missing text} & \key{--} & 空参数或 NA 的显示标签。\\
\key{number format} & \key{fixed} & fixed、grouped 或 scientific。\\
\key{band} & 空 & 范围内严格递增的两个参考端点。\\
\key{band color} & \key{gold} & 区间背景颜色。\\
\key{band opacity} & \key{0.3} & 区间透明度，范围为 0 到 1。\\
\key{stack colors} & 蓝、青绿、金黄 & 堆叠条分段颜色列表，可循环使用。\\
\bottomrule
\end{tabularx}
\end{center}
\endgroup


\subsection{数据、标签与图例参数}
\begin{center}
\begin{tabularx}{\linewidth}{@{}>{\centering\arraybackslash}p{36mm}>{\centering\arraybackslash}p{23mm}>{\centering\arraybackslash}X@{}}
\toprule
\multicolumn{1}{c}{\textbf{参数}} & \multicolumn{1}{c}{\textbf{默认值}} & \multicolumn{1}{c}{\textbf{含义}}\\
\midrule
\key{stack names} & 空 & 与各段一一对应的名称列表。\\
\key{legend} & \key{false} & 自动在堆叠条下方生成图例。\\
\key{segment labels} & \key{none} & none、value、percent、name、name value、name percent。\\
\key{segment font} & \key{\scriptsize} & 段标签字体。\\
\key{segment text color} & \key{auto} & 自动黑白或显式颜色。\\
\key{label overflow} & \key{auto} & 自动外移或 allow 允许越界。\\
\key{better} & \key{higher} & higher 越大越好，lower 越小越好；只影响条件配色。\\
\key{pattern} & \key{none} & 普通条的纹理名称。\\
\key{pattern color} & \key{black} & 纹理线条或点阵的颜色。\\
\key{stack patterns} & \key{none} & 堆叠段和图例的纹理列表，循环使用。\\
\key{print} & 未应用 & 黑白打印预设。\\
\key{print labels} & \key{false} & 内部标签白底；print 会开启此项。\\
\bottomrule
\end{tabularx}
\end{center}
\key{\gradbarslegend[选项]{名称列表}} 生成可共享的独立图例；
\key{\gradbarscolumn{G}{选项}} 定义统一样式的数值列。

CSV 命令包括 \key{\gradbarsloadcsv}、\key{\gradbarscsvtable}、\key{\gradbarscsvstyle}，\par
以及 \key{\gradbarcsv} 和 \key{\gradbarscsvcell}。
完整参数与范围计算规则见“CSV 数据导入”一节。

\subsection{图形与排版参数}
\begin{center}
\begin{tabularx}{\linewidth}{@{}>{\centering\arraybackslash}p{32mm}>{\centering\arraybackslash}p{22mm}>{\centering\arraybackslash}X@{}}
\toprule
\multicolumn{1}{c}{\textbf{图形参数}} & \multicolumn{1}{c}{\textbf{默认值}} & \multicolumn{1}{c}{\textbf{说明}}\\
\midrule
\key{shape} & \key{bar} & bar 普通条；lollipop 棒棒糖。哑铃与趋势使用专用命令。\\
\key{compare color} & \key{black!20} & 双层条的基准层颜色。\\
\key{compare ratio} & \key{.45} & 当前层高度比，严格在 0 与 1 之间。\\
\key{range point} & 空 & 浮动区间内的点，仅用于 gradrange。\\
\key{marker size} & \key{2pt} & 圆点半径，必须为正，受轨道高度限制。\\
\key{stem width} & \key{.6pt} & 棒棒糖及区间端点线宽，必须为正。\\
\key{outline} & \key{false} & 是否只绘制轮廓。\\
\key{outline width} & \key{.4pt} & 轮廓线宽，必须为正。\\
\key{row padding} & \key{0pt} & 图形上下各增加的留白，必须非负。\\
\key{preset} & 未应用 & paper、report、presentation、outline。\\
\bottomrule
\end{tabularx}
\end{center}


\subsection{图形与评价接口}
\begin{center}
\begin{tabularx}{\linewidth}{@{}>{\raggedright\arraybackslash}>{\centering\arraybackslash}p{43mm}>{\centering\arraybackslash}X@{}}
\toprule
\multicolumn{1}{c}{接口或选项} & \multicolumn{1}{c}{说明}\\\midrule
\key{\graddumbbell[o]{a}{b}} & 基准 a、当前 b；空心与实心圆。\\
\key{compare label} & values（默认）、delta 或 percent；也适用于双层条。\\
\key{stack totals} & net（默认净值）或 separate（正小计 / 负小计）。\\
\key{\gradspark[o]{list}} & 等间距序列；空字段和 NA 为缺失。\\
\key{spark range} & auto（默认行内范围）或 fixed（公共 min/max）。\\
\key{spark points} & extrema（默认）、all、last、none。\\
\key{spark high color} & 默认 gradbarsTeal。\\
\key{spark low color} & 默认 gradbarsOrange。\\
\key{better} & higher、lower、target、interval。\\
\key{quality target} & target 模式的数值目标，默认空，必须提供。\\
\key{quality range} & interval 模式的闭区间，默认空，必须提供。\\
\key{palette} & categorical、sequential、diverging、mono，默认不应用。\\
\key{\gradbarscategory{n}{c}{p}} & 固定名称 n 的颜色 c 和纹理 p，定义遵循作用域。\\
\key{references} & 多条参考记录：数值/颜色/线型/名称；默认空。\\
\key{\gradbullet[o]{v}} & 分级背景与居中的实际值条，需设置 bullet bands。\\
\key{\gradintervals[o]{list}} & 多组区间：名称/下界/点估计/上界/颜色。\\
\bottomrule
\end{tabularx}
\end{center}

\clearpage
\subsection{表格与标注接口}
\begin{center}
\begin{tabularx}{\linewidth}{@{}c>{\centering\arraybackslash}X@{}}
\toprule
\multicolumn{1}{c}{接口} & \multicolumn{1}{c}{含义}\\\midrule
\key{gradtable} & 表格环境：外观、列宽与分页设置。\\
\key{\gradcolumn[o]{id}{title}} & 声明列：type 选择类型，options 接受图形与数字选项。\\
\key{\gradrow{list}} & 按列顺序填写逗号分隔的数据。\\
\key{\gradheader{spans}} & 添加一层表头，用 gradspan 声明连续的合并列。\\
\key{\gradspan{n}{title}} & 在 gradheader 中合并 n 列并设置标题。\\
\key{\gradgroup{title}} & 插入跨列分组标题，与首条数据保持在同页。\\
\key{\gradtablestyle{n}{keys}} & 保存表格选项，通过 style=n 调用。\\
\key{\gradtableschema{n}{body}} & 保存列定义，通过 graduseschema 调用。\\
\key{mark=none|best|both} & 数字、普通条、子弹图与 Heat 列的最佳及次佳值标记。\\
\key{data / source} & CSV 数据集及列字段绑定。\\
\key{scale / include zero} & 固定或自动列尺度，控制是否包含零。\\
\key{type=heat / \gradheat} & 连续色阶单元格及独立命令。\\
\key{width / weight} & 表格总宽度，以及可伸缩列分配剩余空间的权重。\\
\key{long=true} & 单栏正文中的跨页模式，重复完整表头。\\
\key{overflow mark} & none（单图默认）、arrow、slash。\\
\key{marker shape} & circle（默认）、square、diamond、triangle。\\
\key{marker fill} & solid（默认）或 hollow。\\
\key{compare annotation} & none（默认）、arrow、bracket。\\
\key{compare label=percent} & 相对于基准绝对值的百分比变化。\\
\bottomrule
\end{tabularx}
\end{center}

\subsection{便捷命令与精细标签定位}
声明入口为 \key{\gradbarsnewcommand[声明与图形选项]{\命令名}}。
完整教程从第~\ref{sec:command-declarations-cn}~节开始；
声明与定位键的默认值及限制集中在第~\ref{sec:command-reference-cn}~节。
参数形式见第~\ref{sec:command-signatures-cn}~节，独立参数与双值输入见
第~\ref{sec:command-data-cn}~节，配色与尺寸见第~\ref{sec:command-colors-cn}~节。
数字方向见第~\ref{sec:command-positions-cn}~节，轨道与端点参照见
第~\ref{sec:command-reference-points-cn}~节，间距、锚点及偏移见第~\ref{sec:command-offsets-cn}~节。
空白与覆盖顺序见第~\ref{sec:command-scope-cn}~节，表格接入见第~\ref{sec:command-tables-cn}~节。

\subsection{已知边界}
CSV 仅支持逗号分隔的单行字段；数值列可自动计算共享范围。

CSV 不支持字段内换行；旧的 \key{gradbarscsvtable} 不分页，原生 \key{gradtable} 数据源支持 \key{long=true}。趋势线使用等间距观测，不解析日期。
不要把绘图命令直接放入 siunitx 的数值型 S 列或 PDF 书签。
宏包已通过 pdfLaTeX、XeLaTeX、LuaLaTeX 的功能检查（包括 CSV）。本中文手册采用 XeLaTeX，英文手册可用 pdfLaTeX 编译。Tagged PDF 支持见第~\ref{sec:manual-accessibility}~节；文档的整体无障碍合规需要独立验证。
